\noindent\textbf{ZUR NETZAUSGLEICHUNG.}\hfill323
\begin{center}\textsc{Gauss} an \textsc{Gerling}.\quad Göttingen, 2. Mai 1837.\end{center}

.\ .\ .\ .\ .\ .\ Was die Berechnung der geodätischen Messungen betrifft, so
habe ich die Elimination bei meinen eigenen Hauptdreiecken von Göttingen
bis Jever, etwa 40 Dreiecken zwischen 33 Punkten, zwar vollständig und
mit aller Strenge durchgeführt. Trotz der Kunstgriffe, die ich dabei
freilich angewandt habe, war das damals aber eine sehr beschwerliche und
langwierige Arbeit. Allein diese Kunstgriffe Ihnen ausreichend zu erklären,
würde ein kleines Buch erfordern. Um dieses auszuarbeiten, müsste ich mich
erst wieder in die Sache einarbeiten, woran ich jetzt gar nicht denken kann.

\begingroup
\def\letterfourstars{*}
\makeatletter
\renewcommand{\@makefnmark}{\hbox{\normalfont[\textsuperscript{\letterfourstars})]}}
\renewcommand{\@makefntext}[1]{\parindent1em\noindent\hbox to2em{\hss[\textsuperscript{\letterfourstars})\ }#1]}
\makeatother
Bei den später von meinem Sohn, \textsc{Hartmann} und \textsc{Müller} ausgeführten
Dreiecken in Westfalen, im Lüneburgischen, im Harz und in der Wesergegend
sowie bei denen des vergangenen Jahres habe ich die äußerste Strenge
etwas gelockert und bin so vorgegangen, dass ich
\begin{list}{}{\setlength{\leftmargin}{2.4em}\setlength{\labelwidth}{1.5em}\setlength{\labelsep}{.5em}\setlength{\itemsep}{.3em}\setlength{\topsep}{.3em}}
\item[1)] die Winkelbedingungen \textls[120]{allein} berücksichtigt und danach streng ausgeglichen habe. Dann habe ich
\item[2)] zu diesen ausgeglichenen Zahlen die neue Ausgleichung gesucht, die
die Seitenverhältnisse erfordern\footnote{Vgl. S.\ 253.}.
\end{list}
\def\letterfourstars{**}
\par\noindent\hspace*{\parindent}Wenn man sich die Mühe macht, dieses Verfahren wiederholt im Wechsel
anzuwenden, bis sich die Ergebnisse stabilisieren, wie ich es bereits im
Supplementum etc.\ p.\ 28\footnote{Band IV, S.\ 77/79.} empfohlen habe, gelangt man,
wie dort bewiesen ist, zu \textls[120]{genau} denselben Ergebnissen, wie wenn man
alle Bedingungen auf einmal berücksichtigt.
\par\medskip
.\ .\ .\ .\ .\ .
\endgroup
